\documentclass[11pt]{article}

\usepackage[letterpaper,margin=1in]{geometry}
\usepackage[T1]{fontenc}
\usepackage[utf8]{inputenc}
\usepackage{lmodern}
\usepackage{microtype}
\usepackage{amsmath,amssymb,amsthm,mathtools}
\usepackage{booktabs,longtable,array}
\usepackage{enumitem}
\usepackage{xcolor}
\usepackage{hyperref}
\usepackage[nameinlink,capitalise]{cleveref}
\usepackage{listings}
\usepackage{fancyhdr}
\usepackage{url}
\setlength{\emergencystretch}{2em}

\hypersetup{
  colorlinks=true,
  linkcolor=blue!50!black,
  citecolor=green!35!black,
  urlcolor=blue!60!black,
  pdftitle={Dyadic-Comb Quadrature for the Fabius and Rvachev Functions},
  pdfauthor={Research report prepared for Vladimir Reshetnikov}
}

\pagestyle{fancy}
\fancyhf{}
\fancyhead[L]{Dyadic-comb quadrature for Fabius and Rvachev functions}
\fancyhead[R]{\thepage}
\setlength{\headheight}{14pt}

\newtheorem{theorem}{Theorem}[section]
\newtheorem{proposition}[theorem]{Proposition}
\newtheorem{lemma}[theorem]{Lemma}
\newtheorem{corollary}[theorem]{Corollary}
\newtheorem{conjecture}[theorem]{Conjecture}
\theoremstyle{definition}
\newtheorem{definition}[theorem]{Definition}
\newtheorem{remark}[theorem]{Remark}
\newtheorem{example}[theorem]{Example}

\DeclareMathOperator{\sinc}{sinc}
\DeclareMathOperator{\ord}{ord}
\DeclareMathOperator{\sgn}{sgn}
\DeclareMathOperator{\FP}{FP}
\newcommand{\Up}{\operatorname{up}}
\newcommand{\eps}{\varepsilon}
\newcommand{\dd}{\,\mathrm{d}}
\newcommand{\e}{\mathrm{e}}
\newcommand{\ii}{\mathrm{i}}
\newcommand{\Z}{\mathbb{Z}}
\newcommand{\N}{\mathbb{N}}
\newcommand{\R}{\mathbb{R}}
\newcommand{\C}{\mathbb{C}}
\newcommand{\Q}{\mathbb{Q}}
\newcommand{\Mell}{\mathcal{M}}
\newcommand{\fall}[2]{#1^{\underline{#2}}}
\newcommand{\rise}[2]{#1^{\overline{#2}}}
\newcommand{\bigO}{\mathcal{O}}

\lstdefinestyle{pythonstyle}{
  language=Python,
  basicstyle=\ttfamily\small,
  keywordstyle=\color{blue!55!black},
  commentstyle=\color{green!35!black},
  stringstyle=\color{red!45!black},
  showstringspaces=false,
  breaklines=true,
  frame=single,
  columns=fullflexible
}

\title{\bfseries Dyadic-Comb Quadrature for the Fabius and Rvachev Functions\\[0.4em]
\large Exact Polynomial Moments, Fractional Mellin Corrections, and Iterated Sums}
\author{Research report prepared for Vladimir Reshetnikov}
\date{28 August 2026}

\begin{document}
\maketitle

\begin{abstract}
Let $F$ be the Fabius function, let $\Up$ be Rvachev's compactly supported
$C^\infty$ function, and put $h=2^{-m}$.  This report studies weighted sums of
$F$ and $\Up$ on the dyadic comb $h\Z$, with particular emphasis on monomial
weights.  The main result is an exact shifted quadrature theorem:
\[
 h\sum_{k\in\Z}\bigl(h(k+\theta)\bigr)^p
 \Up\bigl(h(k+\theta)\bigr)
 =\int_{-1}^{1}x^p\Up(x)\dd x
 \qquad(0\le p\le m),
\]
valid for every real shift $\theta$.  Its mechanism is the increasing
multiplicity of the integer zeros of the infinite sinc product
$\widehat{\Up}$.  We derive the complete finite $2$-adic alias decomposition,
a closed first-failure series, phase-dependent superconvergence, and a
Bell--Bernoulli formula for every higher alias jet.

For the one-sided Fabius comb we obtain a second exact stabilization law.  If
$M=2^m$, $h=M^{-1}$, $p\in\N_0$, and
$\tau(0)=0$, $\tau(p)=2\lceil p/2\rceil$ for $p\ge1$, then
\[
 h\sum_{a=0}^{M}(ah)^pF(ah)
 =\frac{h^{p+1}B_{p+1}(M+1)-d_{p+1}}{p+1}
 \qquad(m\ge\tau(p)),
\]
where $d_s$ is the Mellin moment of the Fabius probability density.  An
all-depth Hurwitz-zeta--Thue--Morse formula is also proved.  For fractional
powers, the punctured Rvachev comb has the universal singular correction
$2\zeta(-\alpha)h^{\alpha+1}$, with a shifted version governed by two Hurwitz
zeta values.  Negative powers are treated by Mellin continuation and a
logarithmic transition at $\alpha=-1$.  Finally, the $n$-fold discrete prefix
sum is reduced to a binomial Cauchy kernel and, at full support, to a finite
combination of exact moments.  The report includes exact-arithmetic and
high-precision numerical experiments and records conjectures and directions
for formalization, $q$-filters, inverse-Fabius sums, and Lambert-$W$ remainder
asymptotics.
\end{abstract}

\tableofcontents

\section{Executive summary of the new results}

The following statements are proved in this report.  ``New'' below means new
relative to the live ProveIt exposition and its consolidated frontier volume
at the time of the audit; no universal priority claim is made without a
broader literature search.

\begin{enumerate}[leftmargin=2.2em]
\item \textbf{Shift-independent exact quadrature for $\Up$.}
At depth $m$, every translate of the dyadic comb integrates exactly every
polynomial of degree at most $m$ against $\Up$.

\item \textbf{A finite $2$-adic alias hierarchy.}
For a monomial of degree $p>m$, only the valuation layers
$v_2(\ell)=0,1,\dots,p-m-1$ can contribute to the Poisson error.  Thus a
formally infinite Fourier error has only finitely many arithmetic types.

\item \textbf{First-failure and superconvergent phases.}
The degree-$m+1$ error is an explicit odd-frequency series involving
$\widehat{\Up}(q/2)$.  If $m$ is even, the centered and half-shifted combs gain
one extra exact odd degree.  If $m$ is odd, the quarter-shifted combs gain one
extra exact even degree.

\item \textbf{Bell--Bernoulli alias jets.}
Every derivative of $\widehat{\Up}$ at an integer zero is expressed by a
complete exponential Bell polynomial.  The finite sinc block contributes an
explicit Bernoulli-number logarithmic jet.

\item \textbf{Exact Fabius right-comb stabilization.}
For natural $p$, the right Riemann comb of $x^pF(x)$ becomes exactly equal to a
finite Euler--Maclaurin polynomial as soon as $m\ge2\lceil p/2\rceil$.  The
closed form contains one Bernoulli polynomial and one Fabius-law moment.

\item \textbf{All-depth complex-power formula.}
For arbitrary complex $\alpha$, the finite Fabius comb is a finite
Thue--Morse sum of differences of Hurwitz zeta values.  For natural $p$ this
reduces to Bernoulli polynomials; for negative integers it reduces to
harmonic-number combinations.

\item \textbf{Fractional and negative Rvachev powers.}
For the punctured comb of $|x|^\alpha\Up(x)$, the sole algebraic correction is
$2\zeta(-\alpha)h^{\alpha+1}$; the rest is smaller than every power of $h$.
The Mellin transform has only one pole, at $\alpha=-1$, because
$\Up(x)-1$ is flat at the origin.

\item \textbf{Closed $n$-fold iterated sums.}
The $n$-fold prefix sum has the exact binomial kernel
$\binom{N-j+n-1}{n-1}$.  At the right edge of the support, polynomial
exactness converts it into a finite moment formula for $\Up$ and a finite
Bernoulli--moment formula for $F$.

\item \textbf{An integer-base extension.}
The same proof works for the compact density whose Fourier transform is
$\prod_{j\ge0}\sinc(\pi\xi/b^j)$: the $b$-adic comb at depth $m$ is exact
through degree $m$ for every integer $b\ge2$.
\end{enumerate}

\section{Repository audit and separation from prior work}
\label{sec:audit}

\subsection{Scope of the audit}

The current documentation tree contains the mathematical glossary, the main
Fabius/Rvachev exposition, the non-elementarity article, the Lean
walkthrough, the consolidated non-formalized frontier volume, a papers
library, and a frozen archive of historical versions.  The live TeX sources
were read directly.  The archive manifest was used to identify historical
versions, byte-identical duplicates, and documents absorbed into the
consolidated frontier volume; mathematically distinct archived themes were
then cross-checked against that canonical consolidation.  This avoids
counting successive drafts as independent results while retaining their
mathematical content.

The audit concentrated on the following existing ingredients:
\begin{itemize}[leftmargin=2em]
\item the relation between $F$ and $\Up$ and the compact-support Fourier
product;
\item exact dyadic point values expressed through signed Thue--Morse sums and
moments;
\item the multiplicity and leading derivative of every integer zero of the
infinite sinc product;
\item Poisson summation and the unweighted partition of unity;
\item moment recurrences, Bernoulli-polynomial identities, and $q$-product
representations;
\item continuous repeated integration and separate work on repeated prefix
sums of the signed Thue--Morse sequence;
\item endpoint and inverse-function asymptotics involving Lambert $W$.
\end{itemize}

Targeted searches in the live main exposition and consolidated frontier
volume found no theorem formulated as a shifted dyadic-comb quadrature rule,
no polynomial moment exactness statement of degree growing with the mesh
depth, and no Fabius right-comb stabilization formula of the form proved
below.  The present report therefore treats those statements as new relative
to the audited repository.  Existing repeated-integration formulas are
continuous Volterra identities; the iterated sums in \cref{sec:iterated} are
discrete and use a different kernel.

\subsection{Canonical repository sources}

The live source tree is at
\begin{center}
\url{https://github.com/VladimirReshetnikov/ProveIt/tree/main/Analysis/FabiusFunction/docs}.
\end{center}
The principal mathematical source is
\begin{center}
\url{https://github.com/VladimirReshetnikov/ProveIt/blob/main/Analysis/FabiusFunction/docs/Fabius_Function_and_Rvachev_Up/Fabius_Function_and_Rvachev_Up.tex}.
\end{center}
The canonical non-formalized frontier source is
\begin{center}
\url{https://github.com/VladimirReshetnikov/ProveIt/blob/main/Analysis/FabiusFunction/docs/non-formalized-research-frontiers/non-formalized-research-frontiers.tex}.
\end{center}
Its README records the consolidation of the former dyadic, $q$-connection,
repeated-integration, convergence, and inverse/saddle notebooks.  The frozen
archive and its provenance manifest are at
\begin{center}
\url{https://github.com/VladimirReshetnikov/ProveIt/tree/main/Analysis/FabiusFunction/docs/archive}.
\end{center}

\begin{remark}[Novelty standard]
The exact identities below are proved from established ingredients, not
inferred from numerical data.  Nevertheless, a claim that they have never
appeared anywhere in the literature would require a larger bibliographic
search than was possible here.  Conjectural statements are explicitly
labeled.
\end{remark}

\section{Notation and established ingredients}
\label{sec:background}

\subsection{Fabius and Rvachev normalizations}

Let $\Up:\R\to\R$ be the even $C^\infty$ function supported on $[-1,1]$ and
normalized by
\[
 \int_{-1}^{1}\Up(x)\dd x=1.
\]
On $[0,1]$, the Fabius function and Rvachev function are related by
\begin{equation}
 \Up(x)=F(1-x),\qquad 0\le x\le1,
 \label{eq:up-F-relation}
\end{equation}
and on $[-1,0]$ by $\Up(x)=F(1+x)$.  Equivalently, the Fabius probability
density on $[0,1]$ is
\begin{equation}
 \rho(x):=F'(x)=2\Up(2x-1).
 \label{eq:fabius-density}
\end{equation}
Both $F$ at $0$ and $1-F$ at $1$ are flat: every derivative vanishes at the
corresponding endpoint.  Consequently $\Up(x)-1$ is flat at $x=0$, while
$\Up$ is flat at $x=\pm1$.

We use the Fourier convention
\[
 \widehat f(\xi)=\int_{\R}f(x)\e^{-2\pi\ii\xi x}\dd x.
\]
The Fourier transform of $\Up$ is the infinite sinc product
\begin{equation}
 \Phi(\xi):=\widehat{\Up}(\xi)
 =\prod_{j=0}^{\infty}\sinc\!\left(\frac{\pi\xi}{2^j}\right),
 \qquad \sinc z:=\frac{\sin z}{z}.
 \label{eq:Phi-product}
\end{equation}
It is even, real on the real axis, and rapidly decreasing together with all
its derivatives.

\subsection{Moments}

Define centered moments
\begin{equation}
 \mu_p:=\int_{-1}^{1}x^p\Up(x)\dd x,
 \qquad
 c_n:=\mu_{2n}.
 \label{eq:mu-c}
\end{equation}
Evenness gives $\mu_{2n+1}=0$.  The rational sequence $c_n$ satisfies
\begin{equation}
 c_0=1,
 \qquad
 (2n+1)(2^{2n}-1)c_n
 =\sum_{j=0}^{n-1}\binom{2n+1}{2j}c_j.
 \label{eq:c-recurrence}
\end{equation}
Let
\begin{equation}
 d_s:=\int_0^1 x^s\rho(x)\dd x
 \label{eq:d-complex}
\end{equation}
for complex $s$.  Since $\rho$ is flat at both endpoints, $d_s$ is entire.
For $r\in\N_0$,
\begin{equation}
 d_r=2^{-r}\sum_{j=0}^{\lfloor r/2\rfloor}\binom r{2j}c_j.
 \label{eq:d-from-c}
\end{equation}
Integration by parts gives the continuous Fabius moment
\begin{equation}
 I_p:=\int_0^1x^pF(x)\dd x=\frac{1-d_{p+1}}{p+1}.
 \label{eq:F-continuous-moment}
\end{equation}
More generally, the right side extends to every complex exponent by a
removable singularity at $p=-1$; see \cref{sec:F-fractional}.

\subsection{Integer zero multiplicities}

For a nonzero integer $n$, write
\[
 d(n):=v_2(|n|)+1,
 \qquad n=\sigma 2^{d(n)-1}q,
 \quad q\ge1\text{ odd},\quad \sigma\in\{\pm1\}.
\]
Exactly $d(n)$ factors in \eqref{eq:Phi-product} vanish at $n$, and the
established integer-zero formula is
\begin{equation}
 \ord_{\xi=n}\Phi=d(n),
 \qquad
 \Phi^{(d(n))}(n)
 =-\frac{d(n)!}{n^{d(n)}}\Phi\!\left(\frac q2\right).
 \label{eq:integer-leading-jet}
\end{equation}
This multiplicity growth is the central engine of the comb identities.

\subsection{Exact dyadic point values}

Put $\eps_r=(-1)^{s_2(r)}$, where $s_2(r)$ is the binary digit sum, and define
\begin{equation}
 P_m(t):=\sum_{j=0}^{\lfloor m/2\rfloor}
 \binom m{2j}c_j t^{m-2j},
 \qquad
 C_m:=\frac{2^{-m(m+1)/2}}{m!}.
 \label{eq:Pm-Cm}
\end{equation}
For $M=2^m$ and $0\le a\le M$, the exact point formula is
\begin{equation}
 F\!\left(\frac aM\right)
 =C_m\sum_{r=0}^{a-1}\eps_r P_m(2a-2r-1).
 \label{eq:F-dyadic-point}
\end{equation}
This formula supplies the finite arithmetic representations in
\cref{sec:finite-arithmetic,sec:F-comb} and the exact numerical experiments.

\section{Shifted dyadic-comb quadrature for \texorpdfstring{$\Up$}{up}}
\label{sec:up-comb}

Let
\[
 M=2^m,\qquad h=M^{-1},\qquad \theta\in\R,
\]
and define the shifted weighted comb
\begin{equation}
 Q_{m,p}(\theta)
 :=h\sum_{k\in\Z}\bigl(h(k+\theta)\bigr)^p
 \Up\bigl(h(k+\theta)\bigr).
 \label{eq:Qmp}
\end{equation}
The sum is finite because $\Up$ is compactly supported.

\subsection{Exact Poisson representation}

\begin{lemma}[Shifted Poisson formula]
For every $p\in\N_0$,
\begin{equation}
 Q_{m,p}(\theta)
 =\mu_p+
 \left(-\frac{1}{2\pi\ii}\right)^p
 \sum_{\ell\in\Z\setminus\{0\}}
 \e^{2\pi\ii\ell\theta}\Phi^{(p)}(M\ell).
 \label{eq:Q-Poisson}
\end{equation}
The series is absolutely convergent.
\end{lemma}

\begin{proof}
Apply shifted Poisson summation to $f_p(x)=x^p\Up(x)$:
\[
 h\sum_k f_p(h(k+\theta))
 =\sum_{\ell\in\Z}\e^{2\pi\ii\ell\theta}\widehat f_p(M\ell).
\]
Differentiating the Fourier transform gives
\[
 \widehat f_p(\xi)=\left(-\frac{1}{2\pi\ii}\right)^p\Phi^{(p)}(\xi).
\]
The zero-frequency term is $\mu_p$.  Rapid decay of $\Phi^{(p)}$ gives
absolute convergence.
\end{proof}

\subsection{Exact polynomial moments}

\begin{theorem}[Shifted dyadic exactness]
\label{thm:shifted-exactness}
For every $m\in\N_0$, every real shift $\theta$, and every polynomial $P$ of
degree at most $m$,
\begin{equation}
 h\sum_{k\in\Z}P\bigl(h(k+\theta)\bigr)
 \Up\bigl(h(k+\theta)\bigr)
 =\int_{-1}^{1}P(x)\Up(x)\dd x.
 \label{eq:polynomial-exactness}
\end{equation}
In particular,
\begin{equation}
 Q_{m,p}(\theta)=\mu_p
 \qquad(0\le p\le m).
 \label{eq:monomial-exactness}
\end{equation}
\end{theorem}

\begin{proof}
For $\ell\ne0$,
\[
 \ord_{\xi=M\ell}\Phi
 =v_2(M|\ell|)+1=m+v_2(|\ell|)+1\ge m+1.
\]
Hence $\Phi^{(p)}(M\ell)=0$ for every $p\le m$.  The alias sum in
\eqref{eq:Q-Poisson} vanishes term by term.  Linearity gives the polynomial
statement.
\end{proof}

\begin{corollary}[Partition and exact centered moments]
For every $m$ and $\theta$,
\[
 h\sum_k\Up(h(k+\theta))=1.
\]
For the centered comb,
\[
 h\sum_k(kh)^{2r}\Up(kh)=c_r\quad(2r\le m),
 \qquad
 h\sum_k(kh)^{2r+1}\Up(kh)=0.
\]
The second odd identity in fact holds at every depth by symmetry.
\end{corollary}

\begin{remark}[Strang--Fix interpretation]
The theorem is a quadrature analogue of the Strang--Fix principle: zeros of a
Fourier transform at the nonzero dual lattice, with increasing multiplicity,
force polynomial reproduction.  What is special here is that the generating
function itself is fixed while the zero multiplicity sampled by the dyadic
dual lattice grows linearly with $m$.
\end{remark}

\subsection{Finite \texorpdfstring{$2$}{2}-adic alias decomposition}

For $p>m$, infinitely many aliases remain, but only finitely many valuation
layers can be nonzero.

\begin{proposition}[Finite valuation hierarchy]
\label{prop:alias-layers}
For every $p\in\N_0$,
\begin{equation}
 Q_{m,p}(\theta)-\mu_p
 =\left(-\frac{1}{2\pi\ii}\right)^p
 \sum_{v=0}^{p-m-1}
 \sum_{\substack{q\in\Z\\q\ \mathrm{odd}}}
 \e^{2\pi\ii 2^vq\theta}
 \Phi^{(p)}(2^{m+v}q),
 \label{eq:alias-layers}
\end{equation}
where the outer sum is empty when $p\le m$.
\end{proposition}

\begin{proof}
Write every nonzero alias uniquely as $\ell=2^vq$ with $q$ odd.  Its zero
order is $m+v+1$.  The $p$th derivative vanishes whenever $p<m+v+1$, or
equivalently $v>p-m-1$.
\end{proof}

This decomposition is useful computationally: increasing $p$ activates one
new $2$-adic shell at a time.  It also suggests valuation-sensitive
Richardson filters; see \cref{sec:future}.

\subsection{The first failure and special phases}

\begin{theorem}[Degree $m+1$ error]
\label{thm:first-failure}
For $p=m+1$,
\begin{equation}
\begin{split}
 Q_{m,m+1}(\theta)-\mu_{m+1}
 ={}&-\left(-\frac{1}{2\pi\ii}\right)^{m+1}
 \frac{(m+1)!}{2^{m(m+1)}}\\
 &\times
 \sum_{\substack{\ell\in\Z\\\ell\ \mathrm{odd}}}
 \frac{\Phi(|\ell|/2)}{\ell^{m+1}}
 \e^{2\pi\ii\ell\theta}.
\end{split}
\label{eq:first-failure}
\end{equation}
The right side is not identically zero as a function of $\theta$.  Thus the
uniform degree $m$ in \cref{thm:shifted-exactness} is sharp.
\end{theorem}

\begin{proof}
Only odd aliases survive by \cref{prop:alias-layers}.  At $M\ell$, the zero
order is $m+1$, and \eqref{eq:integer-leading-jet} gives
\[
 \Phi^{(m+1)}(M\ell)
 =-\frac{(m+1)!}{(M\ell)^{m+1}}\Phi(\ell/2).
\]
Substitution into \eqref{eq:Q-Poisson} yields \eqref{eq:first-failure}.  Its
Fourier coefficient at frequency $1$ is nonzero because $\Phi(1/2)\ne0$.
\end{proof}

\begin{corollary}[Phase superconvergence]
\label{cor:phase-super}
At the first failing degree $p=m+1$:
\begin{enumerate}[label=(\roman*),leftmargin=2em]
\item if $m$ is even, then
$Q_{m,m+1}(0)=Q_{m,m+1}(1/2)=\mu_{m+1}=0$;
\item if $m$ is odd, then
$Q_{m,m+1}(1/4)=Q_{m,m+1}(3/4)=\mu_{m+1}$.
\end{enumerate}
\end{corollary}

\begin{proof}
If $m$ is even, then $m+1$ is odd; pairing $\ell$ and $-\ell$ cancels at
$\theta=0$ and $1/2$.  If $m$ is odd, then $m+1$ is even; the pair is
proportional to $\cos(2\pi\ell\theta)$, which vanishes for every odd $\ell$ at
the quarter shifts.
\end{proof}

The complete zero set of the first-failure trigonometric series is not known;
this becomes \cref{conj:phase-classification}.

\subsection{All higher integer jets: Bell and Bernoulli structure}

The leading derivative \eqref{eq:integer-leading-jet} is enough for exactness,
but arbitrary $p$ requires higher jets.  They admit a compact factorization.
For $n\ne0$, write $d=d(n)$, $n=\sigma2^{d-1}q$ as above, and define
\begin{equation}
 P_d^{\mathrm{sinc}}(w)
 :=\prod_{j=0}^{d-1}\sinc\!\left(\frac{\pi w}{2^j}\right).
 \label{eq:finite-sinc-block}
\end{equation}
Factoring the $d$ vanishing sinc terms gives the exact local identity
\begin{equation}
 \Phi(n+w)=w^d U_n(w),
 \qquad
 U_n(w)=-\frac{P_d^{\mathrm{sinc}}(w)}{(n+w)^d}
 \Phi\!\left(\frac q2+\frac{\sigma w}{2^d}\right).
 \label{eq:local-factorization}
\end{equation}

Let $Y_r$ denote the complete exponential Bell polynomial and set
\[
 \lambda_s(n):=\left.\frac{\dd^s}{\dd w^s}\log U_n(w)\right|_{w=0}.
\]

\begin{proposition}[Bell-polynomial alias jets]
\label{prop:Bell-jets}
For every $r\ge0$,
\begin{equation}
 \Phi^{(d+r)}(n)
 =\frac{(d+r)!}{r!}U_n(0)
 Y_r\bigl(\lambda_1(n),\dots,\lambda_r(n)\bigr),
 \qquad
 U_n(0)=-\frac{\Phi(q/2)}{n^d}.
 \label{eq:Bell-jets}
\end{equation}
Moreover,
\begin{equation}
\begin{split}
 \lambda_s(n)={}&
 \left(\log P_d^{\mathrm{sinc}}\right)^{(s)}(0)
 +\sigma^s2^{-ds}
 \left(\log\Phi\right)^{(s)}(q/2)\\
 &+(-1)^s\frac{d(s-1)!}{n^s}.
\end{split}
\label{eq:lambda-s}
\end{equation}
The finite sinc contribution vanishes for odd $s$, while for $s=2j$,
\begin{equation}
 \left(\log P_d^{\mathrm{sinc}}\right)^{(2j)}(0)
 =-\frac{2^{2j-1}|B_{2j}|}{j}\pi^{2j}
 \frac{1-2^{-2jd}}{1-2^{-2j}}.
 \label{eq:Bernoulli-log-jet}
\end{equation}
\end{proposition}

\begin{proof}
Differentiate $w^dU_n(w)$ and use
$U_n^{(r)}(0)=U_n(0)Y_r(\lambda_1,\dots,\lambda_r)$.  Formula
\eqref{eq:lambda-s} follows by logarithmically differentiating
\eqref{eq:local-factorization}.  Finally,
\[
 \log\sinc z
 =-\sum_{j\ge1}\frac{2^{2j-1}|B_{2j}|}{j(2j)!}z^{2j},
\]
and summing the geometric progression over the $d$ sinc factors gives
\eqref{eq:Bernoulli-log-jet}.
\end{proof}

Equations \eqref{eq:alias-layers} and \eqref{eq:Bell-jets} provide a fully
algorithmic exact description for every natural degree $p$: only finitely
many valuation layers are active, and every active derivative is a finite
Bell polynomial in explicit local data.

\section{Finite arithmetic formulas for centered and absolute Rvachev combs}
\label{sec:finite-arithmetic}

The Fourier formulas are structurally transparent.  For exact rational
arithmetic at a fixed depth, it is useful to insert the point-value formula
\eqref{eq:F-dyadic-point} directly.

Define the punctured absolute comb
\begin{equation}
 A_{m,p}:=h\sum_{k\in\Z\setminus\{0\}}|kh|^p\Up(kh),
 \qquad p\in\N_0.
 \label{eq:Amp}
\end{equation}
The puncture removes the ambiguous $0^0$ and is also the correct convention
for negative powers.  Since $\Up(k/M)=F(1-|k|/M)$,
\begin{equation}
 A_{m,p}=\frac{2C_m}{M^{p+1}}
 \sum_{r=0}^{M-2}\eps_r
 \sum_{t=0}^{M-r-2}(M-r-1-t)^pP_m(2t+1).
 \label{eq:absolute-finite-double}
\end{equation}
This is an exact finite rational formula for every $m,p\in\N_0$.

To expose Bernoulli polynomials, put $L=M-r-1$.  Expanding both powers and
using Faulhaber's formula gives
\begin{equation}
\begin{split}
 A_{m,p}={}&\frac{2C_m}{M^{p+1}}
 \sum_{r=0}^{M-2}\eps_r
 \sum_{j=0}^{\lfloor m/2\rfloor}\binom m{2j}c_j
 \sum_{u=0}^{p}(-1)^u\binom puL^{p-u}\\
 &\times\sum_{v=0}^{m-2j}\binom{m-2j}{v}2^v
 \frac{B_{u+v+1}(L)-B_{u+v+1}}{u+v+1}.
\end{split}
\label{eq:absolute-Bernoulli}
\end{equation}
For the signed centered comb,
\[
 Q_{m,p}(0)=
 \begin{cases}
 0,&p\text{ odd},\\
 A_{m,p},&p\text{ even and }p>0,\\
 1,&p=0.
 \end{cases}
\]
When $p\le m$, \eqref{eq:absolute-Bernoulli} collapses to $c_{p/2}$ for even
$p$ by \cref{thm:shifted-exactness}.  Thus the formula yields a family of
nontrivial finite Thue--Morse--Bernoulli identities.

\section{The one-sided Fabius dyadic comb}
\label{sec:F-comb}

Define, first for natural $p$,
\begin{equation}
 R_{m,p}:=h\sum_{a=0}^{M}(ah)^pF(ah),
 \qquad M=2^m,\\ h=M^{-1}.
 \label{eq:Rmp}
\end{equation}
For $p=0$ the term at $a=0$ is zero because $F(0)=0$.

\subsection{All-depth formula for complex powers}

For $\alpha\in\C$, define $0^\alpha F(0)$ by its continuous value $0$.  The
flatness of $F$ at the origin makes this legitimate for every $\alpha$.
Put
\begin{equation}
 R_m(\alpha):=h\sum_{a=0}^{M}(ah)^\alpha F(ah).
 \label{eq:Rmalpha}
\end{equation}

\begin{theorem}[Hurwitz-zeta--Thue--Morse formula]
\label{thm:Hurwitz-F-comb}
For every $m\in\N_0$ and $\alpha\in\C$,
\begin{equation}
\begin{split}
 R_m(\alpha)
 ={}&\frac{C_m}{M^{\alpha+1}}
 \sum_{r=0}^{M-1}\eps_r
 \sum_{j=0}^{\lfloor m/2\rfloor}\binom m{2j}c_j
 \sum_{\ell=0}^{m-2j}\binom{m-2j}{\ell}2^\ell
 (-2r-1)^{m-2j-\ell}\\
 &\times\Bigl[
 \zeta(-\alpha-\ell,r+1)
 -\zeta(-\alpha-\ell,M+1)
 \Bigr].
\end{split}
\label{eq:Hurwitz-F-comb}
\end{equation}
The difference of Hurwitz zeta functions is entire in $\alpha$; any apparent
pole cancels.
\end{theorem}

\begin{proof}
Insert \eqref{eq:F-dyadic-point} into \eqref{eq:Rmalpha}, interchange the
finite sums, and expand
\[
 (2a-2r-1)^{m-2j}
 =\sum_{\ell=0}^{m-2j}\binom{m-2j}{\ell}
 (2a)^\ell(-2r-1)^{m-2j-\ell}.
\]
The remaining finite power sum is
\[
 \sum_{a=r+1}^{M}a^{\alpha+\ell}
 =\zeta(-\alpha-\ell,r+1)-\zeta(-\alpha-\ell,M+1).
\]
\end{proof}

For $\alpha=p\in\N_0$, use
$\zeta(-n,x)=-B_{n+1}(x)/(n+1)$ to obtain a fully rational
Bernoulli--Thue--Morse formula at every depth.  For negative integer
$\alpha$, the same expression becomes a finite combination of generalized
harmonic numbers.  Formula \eqref{eq:Hurwitz-F-comb} is therefore the broadest
closed form in this report: it includes natural, fractional, and negative
powers without changing the finite outer arithmetic.

\subsection{Exact stabilization for natural powers}

Set
\begin{equation}
 \tau(0):=0,
 \qquad
 \tau(p):=2\left\lceil\frac p2\right\rceil
 =\begin{cases}p,&p\text{ even},\\p+1,&p\text{ odd},\end{cases}
 \quad p\ge1.
 \label{eq:tau}
\end{equation}

\begin{theorem}[Exact Fabius right-comb formula]
\label{thm:F-stabilization}
For $p\in\N_0$ and $m\ge\tau(p)$,
\begin{equation}
 \boxed{
 R_{m,p}
 =\frac{h^{p+1}B_{p+1}(M+1)-d_{p+1}}{p+1}.}
 \label{eq:F-stabilized-Bernoulli}
\end{equation}
Equivalently,
\begin{equation}
 R_{m,p}
 =I_p+\frac h2
 +\sum_{r=1}^{\lfloor(p+1)/2\rfloor}
 \frac{B_{2r}}{(2r)!}\fall p{2r-1}h^{2r}.
 \label{eq:F-stabilized-EM}
\end{equation}
Thus the Euler--Maclaurin expansion terminates and becomes exact at the stated
dyadic depth.
\end{theorem}

\begin{proof}
Let $f_p(x)=x^pF(x)$ on $[0,1]$ and define the composite trapezoidal sum
\[
 T_{m,p}:=h\left(\frac{f_p(0)}2+
 \sum_{a=1}^{M-1}f_p(ah)+\frac{f_p(1)}2\right).
\]
Since $f_p(0)=0$ and $f_p(1)=1$,
\begin{equation}
 R_{m,p}=T_{m,p}+\frac h2.
 \label{eq:R-vs-T}
\end{equation}
Poisson summation for the zero extension of $f_p$, with the standard
half-value convention at the endpoints, gives
\[
 T_{m,p}=\sum_{\ell\in\Z}\widehat f_p(M\ell),
 \qquad
 \widehat f_p(\xi)=\int_0^1f_p(x)\e^{-2\pi\ii\xi x}\dd x.
\]
For a nonzero integer $q$, integrate by parts $p+1$ times.  Flatness of $F$ at
$0$ and of $1-F$ at $1$ gives
\begin{equation}
 \widehat f_p(q)
 =-\sum_{s=0}^{p}\frac{\fall p s}{(2\pi\ii q)^{s+1}}
 +\frac{\widehat g_p(q)}{(2\pi\ii q)^{p+1}},
 \qquad
 g_p:=D^{p+1}(x^pF).
 \label{eq:IBP-Fourier}
\end{equation}
Pairing $q$ and $-q$ in the first sum leaves exactly the Bernoulli terms in
\eqref{eq:F-stabilized-EM}.  It remains to show that the last alias sum
vanishes.

Let $\rho=F'$ and $A=\widehat\rho$.  The operator identity
\begin{equation}
 g_p=\prod_{r=2}^{p+1}(xD+r)\rho
 \label{eq:g-operator}
\end{equation}
transforms into
\begin{equation}
 \widehat g_p(\xi)
 =\prod_{s=1}^{p}(s-\xi D_\xi)A(\xi).
 \label{eq:g-hat-operator}
\end{equation}
From \eqref{eq:fabius-density},
\begin{equation}
 A(\xi)=\e^{-\pi\ii\xi}\Phi(\xi/2).
 \label{eq:A-Phi}
\end{equation}
At $q=M\ell$, $\ell\ne0$, the zero order of $A$ is
\begin{equation}
 \ord_{\xi=q}A=m+v_2(|\ell|).
 \label{eq:A-zero-order}
\end{equation}
If $p$ is odd and $m\ge p+1$, this order is greater than $p$, so
$\widehat g_p(q)=0$ term by term.  If $p$ is even and $m>p$, the same is true.
At the boundary case $m=p$, only odd $\ell$ have zero order exactly $p$.
Then every lower derivative vanishes and
\[
 \widehat g_p(q)=(-q)^pA^{(p)}(q).
\]
Moreover, all lower derivatives of $\Phi(\xi/2)$ vanish at $q$, so
\[
 A^{(p)}(q)=\e^{-\pi\ii q}2^{-p}\Phi^{(p)}(q/2).
\]
For even $p$ this is real and even in $q$ (including the case $p=m=0$).
Thus the numerator is even, whereas $(2\pi\ii q)^{p+1}$ changes sign under
$q\mapsto-q$.  The two aliases cancel.
Hence the remainder in \eqref{eq:IBP-Fourier} sums to zero whenever
$m\ge\tau(p)$.

Finally, expanding $B_{p+1}(M+1)$ at $M=1/h$ gives
\[
 \frac{h^{p+1}B_{p+1}(M+1)-1}{p+1}
 =\frac h2+
 \sum_{r=1}^{\lfloor(p+1)/2\rfloor}
 \frac{B_{2r}}{(2r)!}\fall p{2r-1}h^{2r}.
\]
Combining this with \eqref{eq:F-continuous-moment} proves
\eqref{eq:F-stabilized-Bernoulli}.
\end{proof}

\begin{example}
The first cases are
\[
 R_{m,0}=\frac12+\frac h2,
\]
\[
 R_{m,1}=\frac{1-d_2}{2}+\frac h2+\frac{h^2}{12}
 \qquad(m\ge2),
\]
\[
 R_{m,2}=\frac{1-d_3}{3}+\frac h2+\frac{h^2}{6}
 \qquad(m\ge2),
\]
and
\[
 R_{m,3}=\frac{1-d_4}{4}+\frac h2+\frac{h^2}{4}-\frac{h^4}{120}
 \qquad(m\ge4).
\]
\end{example}

\begin{remark}
The proof explains the parity in $\tau(p)$.  Odd $p$ requires all aliases of
$A$ to have order $>p$.  Even $p$ permits the boundary order $p$ because the
remaining Fourier terms cancel in symmetric pairs.
\end{remark}

\subsection{An exact absolute-power corollary for \texorpdfstring{$\Up$}{up}}

The Fabius result yields a stronger exact formula for the absolute centered
comb than polynomial exactness alone.

\begin{theorem}[Absolute integer powers]
\label{thm:absolute-integer}
For $p\in\N_0$ and $m\ge\tau(p)$,
\begin{equation}
 \boxed{
 A_{m,p}
 =\frac{2d_{p+1}}{p+1}
 +2\zeta(-p)h^{p+1}.}
 \label{eq:absolute-integer-closed}
\end{equation}
The first term is the continuous absolute moment
\begin{equation}
 \int_{-1}^{1}|x|^p\Up(x)\dd x
 =\frac{2d_{p+1}}{p+1}.
 \label{eq:absolute-continuous}
\end{equation}
Consequently:
\begin{itemize}[leftmargin=2em]
\item for positive even $p$, the comb is exactly equal to the integral;
\item for odd $p$, the sole error is the exact zeta correction
$2\zeta(-p)h^{p+1}$;
\item for $p=0$, the punctured comb is $1-h$, since $2\zeta(0)=-1$.
\end{itemize}
\end{theorem}

\begin{proof}
For $p\ge1$, use \eqref{eq:up-F-relation} and change $a=M-k$:
\[
 A_{m,p}=2h\sum_{k=1}^{M}(kh)^pF(1-kh)
 =2\sum_{j=0}^{p}(-1)^j\binom pj R_{m,j}.
\]
The same binomial combination of the integrals $I_j$ equals
$2d_{p+1}/(p+1)$.  Insert \eqref{eq:F-stabilized-EM}.  Alternating finite
differences annihilate every falling factorial of degree $<p$; the only
surviving endpoint term occurs when $p$ is odd and has degree $p$.  It equals
$2\zeta(-p)h^{p+1}$ by $\zeta(-p)=-B_{p+1}/(p+1)$.  The case $p=0$ follows
from the partition of unity after deleting the central term $h\Up(0)=h$.
\end{proof}

This theorem is the integer resonance of the fractional-power expansion in
the next section.

\section{Fractional and negative powers}
\label{sec:fractional}

\subsection{Branch conventions}

On a two-sided comb, $x^\alpha$ is ambiguous for noninteger $\alpha$.  We
therefore separate
\[
 |x|^\alpha
 \qquad\text{and}\qquad
 \sgn(x)|x|^\alpha.
\]
A chosen complex branch can be reconstructed.  For example, with the
principal value $(-x)^\alpha=\e^{\pi\ii\alpha}x^\alpha$ for $x>0$,
\begin{equation}
 x^\alpha
 =\frac{1+\e^{\pi\ii\alpha}}2|x|^\alpha
 +\frac{1-\e^{\pi\ii\alpha}}2\sgn(x)|x|^\alpha.
 \label{eq:branch-decomposition}
\end{equation}

\subsection{The Mellin transform of \texorpdfstring{$\Up$}{up}}

For $\Re\alpha>-1$, define
\begin{equation}
 \Mell_{\Up}(\alpha)
 :=\int_{-1}^{1}|x|^\alpha\Up(x)\dd x
 =2\int_0^1x^\alpha\Up(x)\dd x.
 \label{eq:Mellin-up}
\end{equation}
Since $\Up(x)-1$ is flat at $0$,
\begin{equation}
 \Mell_{\Up}(\alpha)
 =\frac{2}{\alpha+1}
 +2\int_0^1x^\alpha(\Up(x)-1)\dd x.
 \label{eq:Mellin-continuation}
\end{equation}
The second term is entire in $\alpha$.  Therefore $\Mell_{\Up}$ extends
meromorphically to $\C$ with exactly one pole, at $\alpha=-1$, of residue
$2$.

Using $\Up(x)=F(1-x)$ and the symmetric Fabius density gives the useful
identity
\begin{equation}
 \Mell_{\Up}(\alpha)=\frac{2d_{\alpha+1}}{\alpha+1},
 \label{eq:Mellin-d}
\end{equation}
first for $\Re\alpha>-1$ and then meromorphically.  Thus the generalized
Fabius moment $d_s$ is the natural analytic continuation of all absolute
Rvachev moments.

\subsection{Universal zeta correction}

For every real $\alpha$, define the punctured comb
\begin{equation}
 A_m(\alpha):=h\sum_{k\in\Z\setminus\{0\}}|kh|^\alpha\Up(kh).
 \label{eq:A-alpha}
\end{equation}
For $\alpha>-1$ this is a quadrature approximation to
$\Mell_{\Up}(\alpha)$.

\begin{lemma}[Cutoff-monomial summation]
\label{lem:cutoff-mellin}
Let $\chi\in C_c^\infty([0,\infty))$ be equal to $1$ near $0$, let
$0<\theta\le1$, and let $\beta>-1$.  Then, for every $N>0$,
\begin{equation}
\begin{split}
 h\sum_{k=0}^{\infty}\bigl(h(k+\theta)\bigr)^\beta
 \chi\bigl(h(k+\theta)\bigr)
 ={}&\int_0^\infty x^\beta\chi(x)\dd x\\
 &+\zeta(-\beta,\theta)h^{\beta+1}+\bigO_{\beta,N,\chi}(h^N).
\end{split}
\label{eq:cutoff-mellin}
\end{equation}
The formula continues meromorphically in $\beta$ if the integral is read as
its Mellin finite part.  The only collision of the two displayed residues is
at $\beta=-1$.
\end{lemma}

\begin{proof}
Put
\[
 \chi^*(s):=\int_0^\infty \chi(x)x^{s-1}\dd x.
\]
For $\Re s>0$, integration by parts gives
\[
 \chi^*(s)=-\frac1s\int_0^\infty\chi'(x)x^s\dd x.
\]
The integral on the right is entire and rapidly decreasing in every vertical
strip, because $\chi'$ is supported in a compact subinterval of
$(0,\infty)$.  Hence $\chi^*$ is meromorphic with a single pole at $s=0$,
of residue $1$, and has rapid vertical decay away from that pole.

For $c>\beta+1$, Mellin inversion and absolute convergence give
\[
 h^{\beta+1}\sum_{k=0}^\infty(k+\theta)^\beta
 \chi\bigl(h(k+\theta)\bigr)
 =\frac1{2\pi\ii}\int_{c-\ii\infty}^{c+\ii\infty}
 \chi^*(s)\zeta(s-\beta,\theta)h^{\beta+1-s}\dd s.
\]
Move the contour to $\Re s=-A$.  The Hurwitz zeta pole at $s=\beta+1$
contributes $\chi^*(\beta+1)=\int x^\beta\chi(x)\dd x$, and the pole of
$\chi^*$ at $s=0$ contributes
$\zeta(-\beta,\theta)h^{\beta+1}$.  There are no other poles because $\chi$
is locally constant at the origin.  Polynomial vertical-strip bounds for the
Hurwitz zeta function and the rapid decay of $\chi^*$ make the shifted
integral $\bigO(h^{\beta+1+A})$.  Taking $A$ arbitrarily large proves
\eqref{eq:cutoff-mellin}.  The same contour identity supplies the stated
meromorphic continuation.
\end{proof}

\begin{theorem}[Fractional absolute-power asymptotic]
\label{thm:fractional-zeta}
Fix real $\alpha>-1$.  For every $N>0$,
\begin{equation}
 \boxed{
 A_m(\alpha)
 =\Mell_{\Up}(\alpha)
 +2\zeta(-\alpha)h^{\alpha+1}
 +\bigO_{\alpha,N}(h^N).}
 \label{eq:fractional-zeta}
\end{equation}
For nonnegative integer $\alpha=p$, the remainder becomes identically zero
once $m\ge\tau(p)$, by \cref{thm:absolute-integer}.
\end{theorem}

\begin{proof}
Choose a cutoff $\chi$ as in \cref{lem:cutoff-mellin}, supported inside the
positive half of the support.  On $[0,\infty)$ write
\[
 x^\alpha\Up(x)=x^\alpha\chi(x)
 +x^\alpha\bigl(\Up(x)-\chi(x)\bigr).
\]
The second summand extends by zero to a $C_c^\infty$ function: flatness of
$\Up-1$ removes the apparent singularity at $0$, and flatness of $\Up$
handles the support endpoint.  Poisson summation therefore makes its comb
error $\bigO(h^N)$ for every $N$.  Apply
\cref{lem:cutoff-mellin} with $\theta=1$ to the first summand.  The positive
half-comb contributes $\zeta(-\alpha)h^{\alpha+1}$, and evenness supplies an
identical negative-half contribution.  Adding the two integral terms gives
$\Mell_{\Up}(\alpha)$.
\end{proof}

The striking feature is the absence of an ordinary infinite Euler--Maclaurin
series: the flatness of $\Up-1$ kills every higher local coefficient.  At
integer powers the zeta value explains the parity phenomenon:
$\zeta(-2r)=0$, while $\zeta(-(2r+1))$ is a nonzero Bernoulli number.

\subsection{Shifted fractional combs}

For $0<\theta<1$, define
\begin{align}
 A_{m,\theta}(\alpha)
 &:=h\sum_{k\in\Z}|h(k+\theta)|^\alpha
 \Up(h(k+\theta)),
 \label{eq:shifted-abs-alpha}\\
 S_{m,\theta}(\alpha)
 &:=h\sum_{k\in\Z}\sgn(k+\theta)|h(k+\theta)|^\alpha
 \Up(h(k+\theta)).
 \label{eq:shifted-signed-alpha}
\end{align}

\begin{theorem}[Shifted Hurwitz corrections]
\label{thm:shifted-Hurwitz}
For fixed real $\alpha>-1$, fixed $0<\theta<1$, and every $N>0$,
\begin{align}
 A_{m,\theta}(\alpha)
 ={}&\Mell_{\Up}(\alpha)
 +h^{\alpha+1}\bigl[
 \zeta(-\alpha,\theta)+\zeta(-\alpha,1-\theta)
 \bigr]
 +\bigO(h^N),
 \label{eq:shifted-abs-Hurwitz}\\
 S_{m,\theta}(\alpha)
 ={}&h^{\alpha+1}\bigl[
 \zeta(-\alpha,\theta)-\zeta(-\alpha,1-\theta)
 \bigr]
 +\bigO(h^N).
 \label{eq:shifted-signed-Hurwitz}
\end{align}
\end{theorem}

\begin{proof}
Split the shifted lattice at the origin.  Its positive distances are
$h(k+\theta)$ for $k\ge0$, while its negative distances are
$h(k+1-\theta)$ for $k\ge0$.  Apply
\cref{lem:cutoff-mellin} with shifts $\theta$ and $1-\theta$ to the two local
monomials, and use Poisson summation for the two smooth flat remainders.  The
absolute comb adds the two singular contributions.  The signed comb
subtracts them, and its continuous integral vanishes because $\Up$ is even.
\end{proof}

The two Hurwitz zeta values are the right- and left-hand singular lattice
contributions.  At natural integer powers, the reflection identity for
Bernoulli polynomials makes these formulas consistent with the exact shifted
quadrature theorem and the special phases in \cref{cor:phase-super}.

\subsection{Negative powers and the logarithmic transition}

For $-1<\alpha<0$, \eqref{eq:fractional-zeta} is an ordinary convergent
quadrature statement.  At and below $-1$, the continuous integral diverges,
but the Mellin continuation gives a canonical renormalization.

\begin{proposition}[Renormalized negative powers]
\label{prop:negative-powers}
For fixed real $\alpha<-1$, define the meromorphic value of
$\Mell_{\Up}(\alpha)$ by \eqref{eq:Mellin-continuation}.  Then, away from
$\alpha=-1$,
\begin{equation}
 A_m(\alpha)-2\zeta(-\alpha)h^{\alpha+1}
 \longrightarrow \Mell_{\Up}(\alpha)
 \qquad(m\to\infty),
 \label{eq:negative-renormalized}
\end{equation}
with a superalgebraic remainder after the singular term is removed.
At $\alpha=-1$,
\begin{equation}
 A_m(-1)
 =2\log M+2\gamma
 +2\int_0^1\frac{\Up(x)-1}{x}\dd x+\bigO_N(h^N)
 \qquad\text{for every }N>0.
 \label{eq:alpha-minus-one}
\end{equation}
\end{proposition}

\begin{proof}
Use the meromorphic continuation in \cref{lem:cutoff-mellin} for the singular
cutoff term and Poisson summation for the flat remainder.  This proves
\eqref{eq:negative-renormalized}, in fact with an $\bigO_{\alpha,N}(h^N)$
remainder for every $N$.

For the transition point write $\alpha=-1+\varepsilon$.  From
\eqref{eq:Mellin-continuation},
\[
 \Mell_{\Up}(-1+\varepsilon)
 =\frac2\varepsilon
 +2\int_0^1\frac{\Up(x)-1}{x}\dd x+O(\varepsilon),
\]
whereas
\[
 2\zeta(1-\varepsilon)h^\varepsilon
 =-\frac2\varepsilon+2\log M+2\gamma+O(\varepsilon).
\]
The poles cancel.  Taking the finite part at $\varepsilon=0$ gives
\eqref{eq:alpha-minus-one}; the contour remainder is again smaller than every
power of $h$.
\end{proof}

Thus:
\begin{itemize}[leftmargin=2em]
\item $-1<\alpha<0$: ordinary convergence after a decaying correction;
\item $\alpha=-1$: logarithmic divergence;
\item $\alpha<-1$: power divergence
$2\zeta(-\alpha)h^{\alpha+1}$ plus a finite-part limit.
\end{itemize}

\subsection{Fractional and negative powers for \texorpdfstring{$F$}{F}}
\label{sec:F-fractional}

The one-sided Fabius problem is qualitatively simpler at the origin because
$F$ is flat there.  Define
\begin{equation}
 I_F(\alpha):=\int_0^1x^\alpha F(x)\dd x.
 \label{eq:IF-alpha}
\end{equation}
This is an entire function of $\alpha$.  Integration by parts gives
\begin{equation}
 I_F(\alpha)=\frac{1-d_{\alpha+1}}{\alpha+1},
 \label{eq:IF-d-alpha}
\end{equation}
with a removable singularity at $\alpha=-1$.  In particular,
\begin{equation}
 I_F(-1)=-d'_0=-\int_0^1\log x\,\rho(x)\dd x.
 \label{eq:IF-minus-one}
\end{equation}
Every negative power is integrable because $F(x)=O(x^N)$ for every $N$ as
$x\downarrow0$.

For the finite comb $R_m(\alpha)$, \cref{thm:Hurwitz-F-comb} is exact at every
depth.  Its large-$m$ Poincare expansion is especially simple.

\begin{theorem}[Fabius Euler--Maclaurin expansion]
\label{thm:F-fractional-EM}
For fixed $\alpha\in\C$ and every $K\ge1$,
\begin{equation}
 R_m(\alpha)
 =I_F(\alpha)+\frac h2
 +\sum_{r=1}^{K-1}\frac{B_{2r}}{(2r)!}
 \fall\alpha{2r-1}h^{2r}
 +\bigO_{\alpha,K}(h^{2K}).
 \label{eq:F-fractional-EM}
\end{equation}
For $\alpha=p\in\N_0$, the falling factorials terminate and
\cref{thm:F-stabilization} shows that the finite series is exactly equal to
the comb once $m\ge\tau(p)$.
\end{theorem}

\begin{proof}
The function $x^\alpha F(x)$, extended by zero at $0$, is $C^\infty$.  At
$x=0$ all derivatives vanish.  At $x=1$, flatness of $1-F$ implies that its
$n$th derivative equals $\fall\alpha n$.  The ordinary Euler--Maclaurin
formula therefore reduces to \eqref{eq:F-fractional-EM}; adding $h/2$ converts
the trapezoidal sum to the endpoint-inclusive right comb.
\end{proof}

\section{Iterated sums of arbitrary order}
\label{sec:iterated}

\subsection{The universal binomial kernel}

Let $a_j$ be a sequence indexed from a fixed lower bound $j_0$, and define the
scaled prefix operator
\[
 (\mathsf P_h a)_N:=h\sum_{j=j_0}^{N}a_j.
\]

\begin{lemma}[Discrete Cauchy kernel]
\label{lem:prefix-kernel}
For every order $n\ge1$,
\begin{equation}
 (\mathsf P_h^n a)_N
 =h^n\sum_{j=j_0}^{N}
 \binom{N-j+n-1}{n-1}a_j.
 \label{eq:prefix-kernel}
\end{equation}
\end{lemma}

\begin{proof}
The formula is immediate for $n=1$.  Summing once more and using the hockey
stick identity
\[
 \sum_{r=j}^{N}\binom{r-j+n-1}{n-1}
 =\binom{N-j+n}{n}
\]
proves the induction step.
\end{proof}

This answers the purely combinatorial part of the iterated-sum problem.  The
Fabius/Rvachev structure enters when the upper endpoint is the right edge of
the support.

\subsection{Full-support \texorpdfstring{$n$}{n}-fold Rvachev sum}

Let
\begin{equation}
 T_{m,p}^{(n)}
 :=h^n\sum_{j=-M}^{M}
 \binom{M-j+n-1}{n-1}(jh)^p\Up(jh).
 \label{eq:T-iter-up}
\end{equation}
Define the polynomial kernel
\begin{equation}
 K_{n,h}(x)
 :=\frac1{(n-1)!}\prod_{s=1}^{n-1}(1-x+sh).
 \label{eq:Knh}
\end{equation}
Then
\begin{equation}
 T_{m,p}^{(n)}
 =h\sum_{j=-M}^{M}(jh)^pK_{n,h}(jh)\Up(jh).
 \label{eq:T-as-comb}
\end{equation}
The weight has degree $p+n-1$, so \cref{thm:shifted-exactness} applies.

\begin{theorem}[Exact $n$-fold Rvachev sum]
\label{thm:iterated-up}
If $m\ge p+n-1$, then
\begin{equation}
 \boxed{
 T_{m,p}^{(n)}
 =\int_{-1}^{1}x^pK_{n,h}(x)\Up(x)\dd x.}
 \label{eq:iter-up-integral}
\end{equation}
Equivalently, if $e_r$ denotes the elementary symmetric polynomial,
\begin{equation}
 T_{m,p}^{(n)}
 =\frac1{(n-1)!}
 \sum_{\ell=0}^{n-1}(-1)^\ell
 e_{n-1-\ell}(1+h,1+2h,\dots,1+(n-1)h)
 \mu_{p+\ell}.
 \label{eq:iter-up-moments}
\end{equation}
\end{theorem}

\begin{proof}
Equation \eqref{eq:T-as-comb} follows from
\[
 h^{n-1}\binom{M-j+n-1}{n-1}
 =\frac1{(n-1)!}\prod_{s=1}^{n-1}(1-jh+sh).
\]
Apply polynomial exactness, then expand the product in powers of $x$.
\end{proof}

Because odd moments vanish, \eqref{eq:iter-up-moments} usually contains only
about half of its nominal terms.  It is a complete closed form in rational
moments $c_r$.

\subsection{Full-support \texorpdfstring{$n$}{n}-fold Fabius sum}

Define
\begin{equation}
 S_{m,p}^{(n)}
 :=h^n\sum_{j=0}^{M}
 \binom{M-j+n-1}{n-1}(jh)^pF(jh).
 \label{eq:S-iter-F}
\end{equation}
As before,
\[
 S_{m,p}^{(n)}
 =h\sum_{j=0}^{M}(jh)^pK_{n,h}(jh)F(jh).
\]

\begin{theorem}[Exact $n$-fold Fabius sum]
\label{thm:iterated-F}
If
\begin{equation}
 m\ge\tau(p+n-1)=2\left\lceil\frac{p+n-1}{2}\right\rceil,
 \label{eq:iter-F-threshold}
\end{equation}
then
\begin{equation}
\boxed{
\begin{aligned}
 S_{m,p}^{(n)}={}&\frac1{(n-1)!}
 \sum_{\ell=0}^{n-1}(-1)^\ell
 e_{n-1-\ell}(1+h,\dots,1+(n-1)h)\\
 &\times
 \frac{h^{p+\ell+1}B_{p+\ell+1}(M+1)-d_{p+\ell+1}}
 {p+\ell+1}.
\end{aligned}}
\label{eq:iter-F-closed}
\end{equation}
\end{theorem}

\begin{proof}
Expand $K_{n,h}$ and apply \cref{thm:F-stabilization} to each power
$p+\ell$.  The threshold in \eqref{eq:iter-F-threshold} is the maximum of
$\tau(p+\ell)$ for $0\le\ell\le n-1$.
\end{proof}

\subsection{Arbitrary upper endpoints and noninteger weights}

At an intermediate dyadic endpoint $N=a<M$, \cref{lem:prefix-kernel} remains
an exact answer:
\[
 h^n\sum_{j=0}^{a}\binom{a-j+n-1}{n-1}(jh)^pF(jh).
\]
For natural $p$, inserting \eqref{eq:F-dyadic-point} and expanding the
polynomial kernel reduces the result to finite Bernoulli-polynomial sums.  For
complex $\alpha$, the same manipulation produces finite differences of
Hurwitz zeta functions, exactly as in \cref{thm:Hurwitz-F-comb}.

For $\Up$ with fractional or negative absolute powers, the kernel
$K_{n,h}(x)$ has a Taylor series at the singular point $x=0$.  Applying the
shifted singular Euler--Maclaurin theorem term by term yields corrections
\[
 \sum_{r\ge0} \kappa_r\zeta(-\alpha-r)h^{\alpha+r+1},
\]
where $\kappa_r$ is the $r$th Taylor coefficient of the even part of
$K_{n,h}$.  Unlike the unweighted case, more than one zeta term can occur
because the iterated kernel is not locally constant.

\section{Integer-base extension and \texorpdfstring{$q$}{q}-connections}
\label{sec:base-b}

For an integer $b\ge2$, consider
\begin{equation}
 \Phi_b(\xi):=\prod_{j=0}^{\infty}
 \sinc\!\left(\frac{\pi\xi}{b^j}\right).
 \label{eq:Phi-b}
\end{equation}
It is the Fourier transform of a compactly supported $C^\infty$ probability
density $\Up_b$, obtained as an infinite convolution of centered uniform
laws at geometrically decreasing scales.

\begin{theorem}[$b$-adic comb exactness]
\label{thm:b-adic}
Let $h=b^{-m}$.  For every real shift $\theta$ and every polynomial $P$ of
degree at most $m$,
\begin{equation}
 h\sum_{k\in\Z}P(h(k+\theta))\Up_b(h(k+\theta))
 =\int_{\R}P(x)\Up_b(x)\dd x.
 \label{eq:b-adic-exactness}
\end{equation}
\end{theorem}

\begin{proof}
At every nonzero dual-lattice point $b^m\ell$, the factors with
$j=0,1,\dots,m$ all vanish.  Thus $\Phi_b$ has zero order at least $m+1$.
Shifted Poisson summation repeats the proof of
\cref{thm:shifted-exactness}.
\end{proof}

This extension clarifies which part of the phenomenon is genuinely dyadic:
the exactness requires an integer geometric base so that the Fourier zeros
align with the dual lattice.  The finer valuation and Thue--Morse structure
is special to $b=2$.

The $q$-connections appear at several levels:
\begin{itemize}[leftmargin=2em]
\item the finite sinc block in \eqref{eq:finite-sinc-block} is a geometric
product whose logarithmic jets are finite $q$-geometric sums;
\item the moment generating function
\[
 \int\e^{tx}\Up(x)\dd x
 =\prod_{j=0}^{\infty}\frac{\sinh(t/2^{j+1})}{t/2^{j+1}}
\]
produces Bernoulli and Bell polynomials by logarithmic expansion;
\item valuation-layer cancellation suggests Gaussian-binomial or
$q$-difference filters that annihilate selected alias shells;
\item for a nonintegral scale ratio $q^{-1}$, exact lattice alignment is lost,
but finite linear combinations of nearby scales may restore prescribed
orders of cancellation.
\end{itemize}

\section{Numerical experiments}
\label{sec:numerics}

All rational tests were performed with Python's \texttt{Fraction}; no
floating-point tolerance was used.  The accompanying file
\texttt{fabius\_dyadic\_comb\_experiments.py} contains the exact point-value
formula, moment recurrences, Bernoulli polynomials, all-depth comb formula,
shifted tests, and iterated sums.  The fractional tests use \texttt{mpmath}
with 70 decimal digits.

\subsection{Moments}

The first centered moments are
\begin{center}
\begin{tabular}{c|c}
\toprule
$n$ & $c_n=\int x^{2n}\Up(x)\dd x$\\
\midrule
0 & $1$\\
1 & $1/9$\\
2 & $19/675$\\
3 & $583/59535$\\
4 & $132809/32531625$\\
5 & $46840699/24405225075$\\
\bottomrule
\end{tabular}
\end{center}
The first Fabius-law moments are
\[
 d_0=1,\quad d_1=\frac12,\quad d_2=\frac5{18},\quad
 d_3=\frac16,\quad d_4=\frac{143}{1350},\quad
 d_5=\frac{19}{270}.
\]

\subsection{Observed stabilization thresholds}

For $0\le p\le10$, the first depth at which
\eqref{eq:F-stabilized-Bernoulli} holds exactly is:
\begin{center}
\begin{tabular}{c|rrrrrrrrrrr}
\toprule
$p$&0&1&2&3&4&5&6&7&8&9&10\\
\midrule
first exact $m$&0&2&2&4&4&6&6&8&8&10&10\\
$\tau(p)$&0&2&2&4&4&6&6&8&8&10&10\\
\bottomrule
\end{tabular}
\end{center}
This supports the optimality conjecture in \cref{conj:threshold-optimal}.

\subsection{Shifted first-failure errors}

The following exact errors illustrate \cref{cor:phase-super}.
\begin{center}
\begin{tabular}{c|c|c|c}
\toprule
$m$ & $p=m+1$ & shift $\theta$ & $Q_{m,p}(\theta)-\mu_p$\\
\midrule
1&2&$0$&$1/72$\\
1&2&$1/4$&$0$\\
2&3&$0$&$0$\\
2&3&$1/4$&$-4643/11059200$\\
3&4&$0$&$-23/5529600$\\
3&4&$1/4$&$0$\\
4&5&$0$&$0$\\
4&5&$1/4$&$4966069/383551316951040$\\
\bottomrule
\end{tabular}
\end{center}
Every tested shift was exactly correct for all degrees $0\le p\le m$.

\subsection{Fractional corrected sequences}

The table displays
\[
 A_m(\alpha)-2\zeta(-\alpha)h^{\alpha+1}.
\]
Rapid stabilization is visible even at modest depths.
\begin{center}
\begin{tabular}{c|c|c|c}
\toprule
$m$ & $\alpha=1/2$ & $\alpha=3/2$ & $\alpha=-1/2$\\
\midrule
5  & 0.4913969866635627 & 0.1707208835823527 & 2.7852273786291746\\
6  & 0.4913969866528344 & 0.1707208835839721 & 2.7852273785097793\\
7  & 0.4913969866527990 & 0.1707208835839732 & 2.7852273785136532\\
8  & 0.4913969866527990 & 0.1707208835839732 & 2.7852273785136552\\
9  & 0.4913969866527990 & 0.1707208835839732 & 2.7852273785136552\\
10 & 0.4913969866527990 & 0.1707208835839732 & 2.7852273785136552\\
\bottomrule
\end{tabular}
\end{center}
For $\alpha=5/2$, the corrected value stabilizes to
$0.0754523281224172153620023427\ldots$.

\subsection{Reproduction instructions}

From the archive directory, run
\begin{lstlisting}[style=pythonstyle]
python fabius_dyadic_comb_experiments.py > numerical_results.txt
\end{lstlisting}
The script has no mandatory third-party dependency for exact tests.  If
\texttt{mpmath} is unavailable, use
\begin{lstlisting}[style=pythonstyle]
python fabius_dyadic_comb_experiments.py --skip-fractional
\end{lstlisting}

\section{Conjectures and frontier directions}
\label{sec:future}

\begin{conjecture}[Optimal Fabius threshold]
\label{conj:threshold-optimal}
For every $p\ge1$, the condition $m\ge\tau(p)$ in
\cref{thm:F-stabilization} is best possible: the stabilized formula fails at
$m=\tau(p)-1$.
\end{conjecture}

This is verified exactly for $1\le p\le10$.  A proof should reduce the first
nonvanishing Fabius alias layer to a signed odd-frequency series and establish
that its value cannot cancel.

\begin{conjecture}[Classification of first-failure phases]
\label{conj:phase-classification}
Modulo $1$, the phase zeros forced by parity in \cref{cor:phase-super} are the
only zeros common to every depth of the corresponding parity:
\begin{itemize}[leftmargin=2em]
\item for even $m$, only $\theta=0,1/2$ are universally superconvergent at
degree $m+1$;
\item for odd $m$, only $\theta=1/4,3/4$ are universally superconvergent at
degree $m+1$.
\end{itemize}
Individual depths may have accidental additional zeros.
\end{conjecture}

\subsection{Promising research programs}

\paragraph{1. Valuation-selective extrapolation.}
Because degree $p$ activates only finitely many shells
$v_2(\ell)\le p-m-1$, combine depths $m,m+1,\dots,m+r$ with coefficients that
annihilate the first $r$ shells.  The coefficients should have a natural
$q$-binomial description with $q=2^{-1}$ or $2^{-p}$.  This may yield exact
finite filters rather than merely asymptotic Richardson extrapolation.

\paragraph{2. Denominator and $2$-adic valuation laws.}
Equations \eqref{eq:absolute-Bernoulli},
\eqref{eq:F-stabilized-Bernoulli}, and \eqref{eq:iter-F-closed} invite exact
valuation analysis.  Determine the reduced denominator of $R_{m,p}$ and of
the iterated sums, and isolate the contributions of Bernoulli denominators,
$m!$, and the Fabius moment denominators.  Von Staudt--Clausen should interact
nontrivially with the dyadic cancellations.

\paragraph{3. Complete alias-kernel sign theory.}
A sign or total-positivity theorem for the odd-frequency coefficients
$\Phi((2r+1)/2)$ would settle threshold optimality and phase uniqueness.  The
cosine expansion of $\Up$ and the signed Thue--Morse product may provide two
complementary approaches.

\paragraph{4. Fractional transseries and Lambert $W$.}
After removing the zeta singular term, the remainder is superalgebraic.  The
repo's endpoint/small-argument asymptotics suggest that its logarithm is
controlled by the same saddle geometry and Lambert-$W$ scales as the Fourier
decay of $\Phi$.  Derive an explicit beyond-all-orders expansion for
\[
 A_m(\alpha)-\Mell_{\Up}(\alpha)-2\zeta(-\alpha)h^{\alpha+1}
\]
uniformly in $\alpha$.

\paragraph{5. Inverse-Fabius combs.}
Study
\[
 h\sum_{k}x_k^pF^{-1}(x_k),\qquad x_k=h(k+\theta),
\]
where the inverse has logarithmic/Lambert-$W$ endpoint behavior.  Singular
Euler--Maclaurin should produce non-power corrections, and the exact inverse
dyadic formulas in the repository may permit arithmetic special cases.

\paragraph{6. Nonintegral geometric ratios.}
For
$\Phi_q(\xi)=\prod_{j\ge0}\sinc(\pi q^j\xi)$ with $q^{-1}\notin\N$, exact
lattice alignment is lost.  Search for finite multiscale filters whose
Fourier symbols vanish to prescribed order at the relevant approximate
aliases.  Gaussian binomial coefficients are a natural candidate.

\paragraph{7. Formalization.}
The proof of \cref{thm:shifted-exactness} requires a Schwartz-class Poisson
theorem, Fourier differentiation, and the already formalized integer zero
multiplicity.  The Fabius stabilization theorem additionally needs a
piecewise-smooth Poisson/trapezoidal lemma and the operator identity
\eqref{eq:g-hat-operator}.  These are modular Lean targets and would fit the
repository's promotion policy.

\paragraph{8. Multidimensional tensor and radial variants.}
Tensor products immediately give exact cubature on dyadic lattices with
coordinatewise degree bounds.  More interesting is a radial construction
whose Fourier transform inherits valuation-dependent zero multiplicities on
spherical lattice shells.

\section{Conclusion}

Dyadic point arithmetic and dyadic-comb quadrature are two faces of the same
Fourier structure.  The point formulas use signed Thue--Morse cancellation in
physical space; the comb formulas use increasing zero multiplicity in
frequency space.  Their combination produces exact results that are stronger
than ordinary numerical quadrature:
\begin{itemize}[leftmargin=2em]
\item arbitrary shifted combs reproduce all $\Up$-moments through degree $m$;
\item natural Fabius combs stabilize to a finite Bernoulli expression at a
sharp-looking parity threshold;
\item fractional powers have a universal zeta correction and no further
algebraic tail;
\item negative powers admit a canonical Mellin finite part;
\item repeated summation closes in a finite moment basis.
\end{itemize}
The resulting framework connects Poisson summation, $2$-adic zero geometry,
Strang--Fix reproduction, Bernoulli and Bell polynomials, Hurwitz zeta values,
Thue--Morse arithmetic, and the Lambert-$W$ asymptotic frontier.  It also
provides several concrete formalization and computation targets.

\appendix

\section{Proof details for the Bernoulli boundary sum}

Starting from \eqref{eq:IBP-Fourier}, the boundary contribution to
$T_{m,p}-I_p$ is
\[
 -\sum_{s=0}^{p}\frac{\fall p s}{(2\pi\ii M)^{s+1}}
 \sum_{\ell\ne0}\ell^{-(s+1)}.
\]
Odd powers of $\ell^{-1}$ cancel.  Put $s=2r-1$.  Then
\[
 -\frac{2\zeta(2r)}{(2\pi\ii)^{2r}}M^{-2r}
 \fall p{2r-1}
 =\frac{B_{2r}}{(2r)!}\fall p{2r-1}h^{2r},
\]
using
\[
 \zeta(2r)=(-1)^{r+1}\frac{B_{2r}(2\pi)^{2r}}{2(2r)!}.
\]
This gives precisely \eqref{eq:F-stabilized-EM} before the spectral remainder
is considered.

\section{Low-order exact identities}

For convenient reference:
\begin{align*}
 \mu_0&=1,&
 \mu_1&=0,&
 \mu_2&=\frac19,&
 \mu_3&=0,&
 \mu_4&=\frac{19}{675},\\
 \Mell_{\Up}(0)&=1,&
 \Mell_{\Up}(1)&=\frac5{18},&
 \Mell_{\Up}(2)&=\frac19,&
 \Mell_{\Up}(3)&=\frac{143}{2700}.&&
\end{align*}
For the punctured absolute comb, once the threshold is reached,
\begin{align*}
 A_{m,0}&=1-h,\\
 A_{m,1}&=\frac5{18}-\frac{h^2}{6},\\
 A_{m,2}&=\frac19,\\
 A_{m,3}&=\frac{143}{2700}+\frac{h^4}{60},\\
 A_{m,4}&=\frac{19}{675}.
\end{align*}
Here $2\zeta(-1)=-1/6$ and $2\zeta(-3)=1/60$.

For the first iterated Rvachev orders,
\begin{align*}
 T_{m,p}^{(1)}&=\mu_p,\qquad m\ge p,\\
 T_{m,p}^{(2)}&=(1+h)\mu_p-\mu_{p+1},\qquad m\ge p+1,\\
 T_{m,p}^{(3)}&=\frac12\left((1+h)(1+2h)\mu_p
 -(2+3h)\mu_{p+1}+\mu_{p+2}\right),\\
 &\hspace{9.4cm}m\ge p+2.
\end{align*}

\section{Source audit table}

\begin{longtable}{>{\raggedright\arraybackslash}p{0.36\textwidth}>{\raggedright\arraybackslash}p{0.57\textwidth}}
\toprule
Source group & Role in this report\\
\midrule
\endfirsthead
\toprule
Source group & Role in this report\\
\midrule
\endhead
\path{FabiusFunction_Mathematical_Glossary.tex}
& Normalizations, notation, and cross-reference map.\\
\path{Fabius_Function_and_Rvachev_Up.tex}
& Fourier product, moments, Poisson summation, integer zero multiplicities,
exact dyadic point values, Bernoulli and $q$ identities.\\
\path{Non_Elementarity_of_the_Fabius_Function.tex}
& Analytic/flatness context; no comb formula duplicated.\\
\path{fabius_lean_walkthrough.tex}
& Formalization boundary and candidate theorem dependencies.\\
\texttt{non-formalized-}\newline\texttt{research-frontiers.tex}
& Consolidated dyadic, $q$, repeated-integration, convergence, and
inverse/saddle research; used to avoid duplication and identify open proof
obligations.\\
Archive manifest and frozen historical versions
& Provenance and identification of earlier drafts, including the separate
$K$-fold signed Thue--Morse summation note.  These are distinguished from the
weighted Fabius/Rvachev comb sums studied here.\\
Papers library
& Arias de Reyna sources and historical background.\\
\bottomrule
\end{longtable}

\section{Implementation notes}

The exact code uses the recurrence \eqref{eq:c-recurrence} and the point
formula \eqref{eq:F-dyadic-point}.  At a fixed depth, all $F(a/2^m)$ are
computed as a convolution of the Thue--Morse signs with the odd samples of
$P_m$.  The implementation is intentionally transparent rather than
asymptotically optimized.  Exact depths through $m\approx10$ are fast enough
for the tables in this report.

The all-depth natural-power check uses
\[
 \sum_{a=r+1}^{M}a^n
 =\frac{B_{n+1}(M+1)-B_{n+1}(r+1)}{n+1},
\]
which is the integer specialization of \eqref{eq:Hurwitz-F-comb}.  Bernoulli
numbers are generated exactly from their triangular recurrence, so the code
does not require a computer algebra system.

\begin{thebibliography}{99}

\bibitem{AriasArithmetic}
J. Arias de Reyna,
\emph{Arithmetic of the Fabius function},
arXiv:1702.06487, version 3, 2017 (manuscript dated 2018).
\url{https://arxiv.org/abs/1702.06487}.

\bibitem{AriasCompact}
J. Arias de Reyna,
\emph{An infinitely differentiable function with compact support: definition
and properties},
Rev. Real Acad. Ciencias Madrid 76 (1982), 21--38;
English version at arXiv:1702.05442.

\bibitem{deBoorRon}
C. de Boor and A. Ron,
\emph{Fourier analysis of the approximation power of principal
shift-invariant spaces},
Constructive Approximation 8 (1992), 427--462.
\url{https://doi.org/10.1007/BF01203462}.

\bibitem{StrangFix}
G. Strang and G. Fix,
\emph{A Fourier analysis of the finite element variational method},
in Constructive Aspects of Functional Analysis, C.I.M.E., 1973,
pp. 793--840.

\bibitem{TrefethenWeideman}
L. N. Trefethen and J. A. C. Weideman,
\emph{The exponentially convergent trapezoidal rule},
SIAM Review 56 (2014), no. 3, 385--458.
\url{https://doi.org/10.1137/130932132}.

\bibitem{Navot1961}
I. Navot,
\emph{An extension of the Euler--Maclaurin summation formula to functions
with a branch singularity},
Journal of Mathematics and Physics 40 (1961), 271--276.

\bibitem{Navot1962}
I. Navot,
\emph{A further extension of the Euler--Maclaurin summation formula},
Journal of Mathematics and Physics 41 (1962), 155--163.
\url{https://doi.org/10.1002/sapm1962411155}.

\bibitem{DLMF}
NIST Digital Library of Mathematical Functions,
chapters on Bernoulli polynomials and the Hurwitz zeta function.
\url{https://dlmf.nist.gov/}.

\bibitem{ProveItMain}
V. Reshetnikov,
\emph{Fabius Function and Rvachev Up}, live TeX exposition in the ProveIt
repository, accessed 28 August 2026.
\url{https://github.com/VladimirReshetnikov/ProveIt/tree/main/Analysis/FabiusFunction/docs}.

\bibitem{ProveItFrontier}
V. Reshetnikov,
\emph{Non-formalized Fabius research frontiers}, consolidated TeX volume in
the ProveIt repository, accessed 28 August 2026.

\end{thebibliography}

\end{document}
